\chapter{Discusión}

Esta discusión se organiza en torno a la \textbf{pregunta principal} de investigación y a las \textbf{preguntas específicas}, contrastando los resultados obtenidos con los criterios y métricas definidos en el método. La pregunta principal planteó: \textit{¿Cuál es la viabilidad técnica y el nivel de usabilidad percibida de una plataforma SaaS modular que integre gestión de pedidos, inventarios, repartidores y atención con inteligencia artificial, al validarse en un entorno controlado en PYMES gastronómicas de Portoviejo?} Los resultados muestran que la plataforma alcanzó los umbrales de éxito definidos (tiempo de registro $\leq$ 3 min, errores $\leq$ 10\,\%, completitud $\geq$ 95\,\%, SUS $\geq$ 68), lo que permite responder afirmativamente en el contexto del piloto simulado: la solución es técnicamente viable y usable bajo los criterios establecidos.

Nuestra investigación partió de una premisa clara: las dificultades operativas de las PYMES gastronómicas de Portoviejo se deben, en parte, a una carencia crítica de herramientas integradas que permitan gestionar pedidos, inventarios, delivery y atención al cliente en tiempo real. Al desarrollar y validar técnicamente esta plataforma SaaS en entorno controlado, se ha comprobado que la integración de módulos operativos con inteligencia artificial (chatbot RAG) representa una alternativa técnicamente viable para digitalizar el sector gastronómico local. Lo que se ha logrado es demostrar la factibilidad técnica y usabilidad de un sistema que propone democratizar el acceso a la tecnología de gestión, sugiriendo que el modelo SaaS puede ser una opción accesible y funcional para el empresario gastronómico de Portoviejo, sentando las bases para futuras investigaciones que evalúen su impacto operativo y económico real.

\section{Interpretación de los hallazgos principales}

\subsection{Pertinencia de los requerimientos y del diseño técnico}

Los requerimientos funcionales identificados incluyen la gestión de pedidos con geolocalización, control de inventarios, administración de repartidores, vista de cocina en móvil, dashboard analítico con KPIs y un chatbot RAG integrado en la plataforma web para atención al cliente. A su vez, los requerimientos no funcionales abarcan usabilidad (SUS $\geq$ 68), rendimiento (tiempo de registro $\leq$ 3 min), fiabilidad (errores $\leq$ 10\,\%), seguridad (autenticación JWT y RBAC) y escalabilidad mediante arquitectura modular. En conjunto, estos criterios resultaron adecuados para representar las necesidades operativas del contexto de estudio y orientar tanto el diseño como la evaluación del artefacto.

La arquitectura adoptada, compuesta por backend en NestJS, aplicación móvil en React Native + Expo, sitio web en Next.js, base de datos relacional en PostgreSQL y autenticación segura, demostró ser apropiada por su mantenibilidad, escalabilidad y costo accesible. La validación técnica confirmó que esta estructura soporta flujos críticos con completitud $\geq$ 95\,\% y estabilidad operativa, lo que refuerza su pertinencia para una solución SaaS dirigida a PYMES gastronómicas.

\subsection{Cumplimiento funcional y desempeño del artefacto}

La plataforma cumple en alto grado los requisitos funcionales críticos especificados, evaluados mediante pruebas de aceptación y trazabilidad de casos. Las métricas obtenidas incluyen tiempos de ejecución adecuados en tareas clave (por ejemplo, 2.5 min para crear un producto y 45 seg para aprobar un repartidor), porcentajes de éxito elevados en tareas guiadas, cobertura aproximada del 57\,\% en módulos base de pruebas unitarias, completitud de integración $\geq$ 95\,\% y cumplimiento de dimensiones de calidad asociadas a funcionalidad, usabilidad, rendimiento, compatibilidad y seguridad. En conjunto, estos resultados permiten sostener que el artefacto alcanzó un nivel satisfactorio de desempeño técnico para el alcance del estudio.

\subsection{Usabilidad e integración del componente de inteligencia artificial}

El módulo de inteligencia artificial se integró mediante un chatbot RAG accesible desde la plataforma web, apoyado en un servicio independiente capaz de procesar consultas y generar respuestas basadas en información recuperada. Esta integración permitió automatizar respuestas a consultas frecuentes sobre pedidos, horarios y servicios, mejorando la rapidez y disponibilidad de la atención al cliente durante las pruebas.

En cuanto a la usabilidad, la plataforma obtuvo un puntaje SUS de 70 en usuarios de prueba, superando el umbral de aceptabilidad de 68. Este resultado indica que el sistema fue percibido como intuitivo y adecuado para tareas operativas del contexto gastronómico. La combinación entre funcionalidad integrada, respuesta técnica estable y apoyo conversacional mediante IA refuerza la idea de que la solución no solo es viable desde el punto de vista técnico, sino también aceptable desde la experiencia de uso.

\section{Viabilidad técnica y significado de los hallazgos}

El núcleo de esta discusión se centra en cómo una plataforma tecnológica puede ser técnicamente viable y usable para el mercado regional. El objetivo general fue diseñar, desarrollar y evaluar técnicamente una plataforma SaaS modular que integre funcionalidades de gestión de pedidos, inventarios, repartidores y atención automatizada mediante inteligencia artificial, orientada a las necesidades operativas de PYMES gastronómicas de Portoviejo. Los resultados de la validación en entorno controlado (\textit{staging}) mediante un escenario piloto simulado permiten afirmar que la solución alcanzó un nivel de operatividad adecuado: la centralización de pedidos, el control de inventarios, la gestión de repartidores y el panel analítico funcionaron de manera integrada en las pruebas, y el chatbot RAG integrado en la plataforma web respondió a consultas frecuentes de los usuarios. Este hallazgo es fundamental para la viabilidad de la propuesta. Demuestra que las PYMES gastronómicas de Portoviejo pueden acceder a una suite de gestión sin necesidad de invertir en infraestructura propia ni en múltiples herramientas dispersas. Además, el despliegue bajo modelo SaaS elimina la fricción de instalaciones complejas y actualizaciones locales, alineándose con la visión moderna de software como servicio que ha demostrado ser efectiva en contextos de pequeña y mediana empresa \parencite{GonzalezMartinez2021}.

La convergencia entre módulos operativos (pedidos, inventarios, repartidores, dashboard) y un componente de inteligencia artificial (chatbot RAG integrado en la plataforma web) responde a la demanda de inmediatez y profesionalización del servicio al cliente, un aspecto crítico en el sector gastronómico y en la competitividad de las ciudades creativas de la gastronomía \parencite{Castells2010}.

Desde una lectura comparativa, los resultados técnicos del piloto se mantienen dentro de umbrales de aceptación reportados para prototipos SaaS validados en contexto de pequeña empresa: completitud de flujos críticos de 96\,\%, estabilidad con 3\,\% de fallos, tiempo de registro de 3 min y SUS de 70. En términos de usabilidad, el puntaje obtenido supera el umbral de aceptabilidad de 68 ampliamente utilizado en la literatura \parencite{Bangor2009, Lewis2018}. En términos de calidad del producto, la verificación por requisitos funcionales, pruebas de integración y trazabilidad de casos es consistente con el enfoque de evaluación por características de calidad planteado por ISO/IEC 25010 \parencite{ISO25010_2011}. Estos resultados no implican generalización estadística, pero sí entregan evidencia técnica comparable para sustentar la viabilidad del artefacto en el contexto de estudio.

\section{Análisis temático de los resultados}

\subsection{Pertinencia de los requerimientos del sistema}

Al iniciar el proyecto, la recolección de requisitos, objetivo general y objetivos específicos se apoyó en revisión de literatura, análisis del dominio y levantamiento exploratorio del problema. El diagnóstico identificó que las PYMES del sector presentan un bajo nivel de digitalización: gestión de pedidos mediante anotaciones o herramientas básicas, inventarios poco sistematizados y atención al cliente dependiente de canales no integrados. Al contrastar este problema con los resultados de la validación técnica, se observó que la plataforma cumplió satisfactoriamente los requisitos funcionales críticos especificados y que la percepción de utilidad por parte de los usuarios de prueba fue favorable. Este fenómeno de ``opacidad operativa''---falta de información centralizada y en tiempo real---es una barrera común reportada en la transformación digital de PYMES en América Latina, donde la ausencia de datos operativos estructurados frena la toma de decisiones \parencite{GonzalezMartinez2021}. La plataforma propuesta demuestra la viabilidad técnica de una solución que contribuye a romper el paradigma de la gestión fragmentada, respondiendo a la necesidad de centralización e inmediatez del negocio gastronómico en Portoviejo.

\subsection{Adecuación de la arquitectura tecnológica}

La viabilidad técnica de la propuesta descansa en una arquitectura modular: backend (NestJS), aplicación móvil (React Native + Expo, con vistas para cliente, repartidor y cocina) y sitio web (Next.js, principalmente administración y también compra como cliente desde la web), con base de datos relacional y autenticación segura (JWT, RBAC). La decisión de adoptar un stack moderno y abierto permitió reducir costos de licenciamiento y facilitar el mantenimiento. Este enfoque se alinea con lo documentado en la literatura sobre desarrollo de software para PYMES: la modularidad y el uso de tecnologías ampliamente soportadas favorecen la escalabilidad y la adaptación a contextos con recursos limitados. El desarrollo iterativo bajo metodología ágil (sprints) permitió ajustar requisitos a partir de la retroalimentación de los participantes de prueba, reforzando la adecuación de la solución al contexto local.

\subsection{Integración de módulos y operatividad de la plataforma}

La implementación de los módulos de gestión de pedidos (con soporte a geolocalización y medios de pago efectivo y transferencia), control de inventarios, vista de cocina en móvil (aceptar o rechazar pedidos, notificar cuando esté listo para que el repartidor recoja), administración de repartidores y dashboard analítico con KPIs permitió a los usuarios de prueba del escenario piloto simulado ejecutar flujos operativos en entorno controlado (\textit{staging}) que antes se realizarían de forma manual o dispersa. La disponibilidad de indicadores en tiempo real (pedidos, ventas, tiempos de entrega) en las pruebas sugiere el potencial de la plataforma para favorecer una toma de decisiones más informada en operación real, aspecto que en la literatura se identifica como uno de los beneficios clave de la analítica de negocio en pequeñas empresas. Los resultados observados en las tareas guiadas (tiempos de registro de pedidos, tasa de errores, completitud de flujos) indican que la plataforma cumple con el propósito de demostrar la viabilidad técnica de optimización de procesos internos y mejora de la calidad del servicio.

\subsection{Aporte del chatbot RAG a la atención al cliente}

La integración del chatbot basado en el modelo RAG (Retrieval-Augmented Generation), accesible desde la plataforma web, evidenció mejoras en la rapidez y la disponibilidad de respuestas a consultas frecuentes (pedidos, horarios, servicios). Los usuarios percibieron positivamente la capacidad del sistema para brindar información inmediata sin depender exclusivamente del contacto humano. Este hallazgo respalda la viabilidad del uso de inteligencia artificial como herramienta de apoyo en la atención al cliente dentro de contextos empresariales de pequeña y mediana escala, en línea con experiencias recientes reportadas para asistentes conversacionales con recuperación contextual \parencite{ChenNiforatosKortuem2025}.

\subsection{Cumplimiento funcional y calidad del sistema}

Esta sección responde a la pregunta específica 4: \textit{¿En qué medida la plataforma desarrollada cumple con los requisitos funcionales especificados y los criterios de calidad de software establecidos, evaluando métricas como tiempo promedio de ejecución de tareas, porcentaje de éxito en tareas guiadas y dimensiones de calidad ISO/IEC 25010?}

La verificación del cumplimiento de requisitos funcionales y de los criterios de calidad de software establecidos se realizó mediante pruebas de integración y validación de flujos críticos en entorno de desarrollo o \textit{staging}. Los resultados permiten afirmar que la plataforma desarrollada cumple con los requisitos funcionales especificados en la fase de levantamiento (gestión de pedidos, inventarios, repartidores, vista cocina, dashboard, chatbot web) y con criterios de calidad adecuados para un prototipo validado en entorno controlado. Las métricas evaluadas incluyen tiempos promedio de tareas (como se detalla en la Tabla 8), porcentajes de éxito en tareas guiadas (100\,\% en la mayoría, con mínimas variaciones en checkout), y cumplimiento de dimensiones ISO/IEC 25010.

\subsection{Usabilidad percibida de la plataforma}

La evaluación mediante la escala SUS (System Usability Scale) y las pruebas de usabilidad en entorno controlado permitieron cuantificar la aceptación del sistema por parte de usuarios de prueba (administradores, cocina y repartidores). Los resultados reflejan que la plataforma fue percibida como intuitiva y alineada con las tareas operativas del contexto gastronómico. En la ingeniería de software se considera que el éxito operativo de un sistema depende de su capacidad para satisfacer las necesidades reales de los usuarios en su entorno de trabajo; la retroalimentación obtenida indica que el diseño centrado en los flujos de pedido, inventario y delivery contribuyó a una experiencia de usuario positiva durante el período de validación técnica. No obstante, se identificaron oportunidades de mejora en la capacitación inicial y en la personalización de algunas funcionalidades, lo que sugiere la necesidad de iteraciones continuas en futuras versiones orientadas a implementación a escala real.

\section{Limitaciones del estudio}

Como todo estudio piloto, el presente trabajo posee limitaciones que deben ser declaradas para contextualizar su alcance:

\oldtextbf{Tamaño de la muestra:} La validación se realizó mediante un escenario piloto simulado representativo del contexto gastronómico de Portoviejo. Aunque este diseño permitió validar la funcionalidad técnica, la usabilidad y el desempeño del sistema en tareas guiadas, no pretende una generalización estadística a toda la población de negocios gastronómicos de la ciudad o de la provincia. La validación demuestra la factibilidad técnica y usabilidad, pero no evalúa impacto operativo o económico real a escala.

\textbf{Alcance funcional:} La plataforma desarrollada corresponde a un prototipo funcional con módulos core (pedidos, inventarios, repartidores, vista cocina, dashboard, chatbot RAG integrado en la plataforma web). Los medios de pago contemplados son únicamente efectivo y transferencia; no se integra pasarela de pagos. Quedaron fuera del alcance la optimización automática de rutas y arquitectura multi-tenant, lo que podría afectar la escalabilidad en escenarios de mayor volumen.

\oldtextbf{Dependencia del contexto piloto:} El rendimiento y la adopción observados están ligados a las condiciones específicas del escenario simulado (volumen de pedidos, perfiles de usuario y conectividad controlada). Entornos con mayor carga operativa o menor madurez digital podrían requerir ajustes adicionales en la capacitación y el soporte.

\textbf{Período de validación:} El período de evaluación técnica de la plataforma en entorno controlado fue acotado (semanas). Una evaluación longitudinal en operación real permitiría observar la sostenibilidad de la adopción y el impacto en indicadores operativos del negocio de manera más estable y representativa.

Los usuarios de prueba sugirieron mejoras para futuras iteraciones: notificaciones push en tiempo real para cocina y repartidores; integración con WhatsApp para confirmaciones automáticas de pedido; filtrado avanzado de pedidos por fechas y estado; modo oscuro en la app móvil; y métodos de pago adicionales (tarjeta, billeteras digitales). Estas recomendaciones se consideran en la sección de limitaciones y trabajo futuro.

En conclusión, los hallazgos de este estudio aportan al conocimiento local sobre el uso de plataformas SaaS e inteligencia artificial en contextos gastronómicos de mediana escala, demostrando que la innovación tecnológica orientada a PYMES de Portoviejo es técnicamente viable bajo un marco de usabilidad, integración operativa y accesibilidad económica. Los resultados de la validación técnica en entorno controlado y escenario simulado sientan las bases para futuras investigaciones orientadas a implementación a escala real y medición de impacto operativo y económico en el sector.
